\section{Stage 4: memory organisations}\label{sec:stage4}

All organisations write to one store that keeps the records, an append-only
operation log, an exact vector index over record texts, and the memory size
$\toklen{M}=\sum_{r\in M}|\tau(\sigma(r))|$, maintained incrementally. The
serialisation $\sigma(r)$ of a record is one line that does not depend on any
other record (names of referenced entities are copied into the record), so the
per-record counts are exact and the budget can be checked after every
operation. \Cref{tab:record-lines} shows the line formats.

\begin{table}[ht]
\centering\small
\caption{Record types and their serialisation $\sigma(r)$.}\label{tab:record-lines}
\begin{tabularx}{\linewidth}{@{}lX@{}}
\toprule
Type & Line \\
\midrule
Entity & \texttt{[E3] entity red mug (container) aka cup: color=red | seen 00:01:20-00:14:00} \\
Event & \texttt{[V12] 00:03:20-00:03:40 event: man pours water into red mug \{E1,E3\}} \\
State change & \texttt{[S4] 00:12:00 change: red mug [E3] contents: water -> empty} \\
Relation & \texttt{[R2] 00:05:00 relation: man [E1] holds drill [E5]} \\
Note (C1, C2) & \texttt{[N7] 00:02:00-00:02:20 note: <free text>} \\
\bottomrule
\end{tabularx}
\end{table}

\subsection{Structured memory (C3)}
For each observation $o_i$ the update $U$ runs five operations in order.

\paragraph{Entity resolution.}
For every observed entity $e$ and every stored entity $x$ of a compatible
category, the similarity is the largest cosine between the embedding of $e$'s
name and those of $x$'s name and aliases, minus a penalty of 0.10 when either
name is generic (``person'', ``tool'', \dots). Candidate pairs with similarity
$\ge\tau_{\mathrm{low}}=0.60$ are processed greedily in decreasing order; an
observed entity takes at most one stored entity and a stored entity is claimed
at most once per chunk (two mugs seen together are two mugs). Pairs above
$\tau_{\mathrm{high}}=0.80$ are linked; pairs in between are put to the LLM,
which sees both descriptions with their attributes and times. Unmatched
entities become new records. Every decision is logged with its similarity, for
the merge audit. Names are compared, not attributes, because attributes
legitimately change (a board is painted); the LLM sees attributes and can use
them.

\paragraph{State changes with debouncing.}
For a linked entity, an observed attribute value that differs from the stored
one opens a pending change. It is committed as a StateChange record, with time
equal to the first observation of the new value and provenance equal to the
observations that confirmed it, only after $k=2$ consecutive observations of
that entity agree. A single noisy description therefore does not create a
false change; the price is that a change is recorded one observation late.
The ablation \emph{no\_debounce} sets $k=1$.

\paragraph{Events, speech and relations.}
Events are inserted with their resolved entity identifiers. An event of the
immediately preceding chunk is extended instead when the new one has the same
entity set, a similar description (cosine $\ge0.85$) \emph{and} no content
word that the stored description lacks (a person stirring for a minute is one
event). The last condition matters: the extended event keeps its original
wording, so ``paints the board blue'' must not be merged into ``paints the
board white''. The speech summary is stored as an event of
modality \emph{speech}. Relations are deduplicated on
(subject, predicate, object).

\paragraph{Budget enforcement.}
While $\toklen{M}>B$: (i)~the oldest uncompressed events outside the three
most recent chunks are grouped into a window of 300~s and replaced by at most
three coarser events written by the LLM, with provenance equal to the union of
the originals (a window holding a single event keeps it verbatim); (ii)~when nothing is left to compress, the record with the
lowest priority is deleted,
\[
\mathrm{prio}(r)=w_{\mathrm{rec}}\,\frac{t(r)-t_{\min}}{t_{\mathrm{now}}-t_{\min}}
+\begin{cases}
w_{\mathrm{sal}}\,\mathrm{sal}(r)+w_{\mathrm{ref}}\,|\mathrm{ents}(r)| & \text{event}\\
w_{\mathrm{sal}}+2w_{\mathrm{ref}} & \text{relation}\\
2w_{\mathrm{sal}}+w_{\mathrm{sc}} & \text{state change}\\
w_{\mathrm{ent}}+w_{\mathrm{ref}}\cdot\#\text{references} & \text{entity}
\end{cases}
\]
with $(w_{\mathrm{sal}},w_{\mathrm{ref}},w_{\mathrm{rec}},w_{\mathrm{sc}},w_{\mathrm{ent}})=(1,0.5,1,2,3)$.
Entities and state changes are never compressed: they are compact and carry
exactly the facts that object-change and cross-segment questions ask for
(the first expectation in \cref{sec:rq}).

\subsection{Baselines}
\begin{description}[nosep]
\item[C1 rolling summary.] After each chunk the LLM rewrites one summary from
  the previous summary and the new observation, under a word limit derived
  from $B$; the result is cut to the longest word prefix whose line fits $B$
  exactly. The answerer receives the whole summary.
\item[C2 flat log.] One note per observation (its text serialisation);
  the oldest note is evicted while $\toklen{M}>B$.
\item[C4 unbounded log.] C2 without eviction: the upper bound for this
  perception. Its final size defines the budgets: $B_k=f_k\cdot\toklen{M_{\mathrm{C4}}}$
  with $f=(0.02,0.05,0.15)$.
\end{description}
The ablations of C3 switch off entity resolution, state changes, compression,
or debouncing, one at a time.

\subsection{Retrieval and answering}\label{sec:retrieval}
For C2--C4, the score of a record is its cosine similarity to the question,
plus 0.25 when it is, or refers to, an entity whose name occurs in the
question, plus $0.20\,w_t$ for temporal cues: \emph{beginning/first} favours
early records, \emph{end/last} late ones, and \emph{after X}/\emph{before X}
favours records just after/before the best match for X with weight
$\exp(-\text{gap}/300\,\mathrm{s})$. The top 40 records are taken in score
order while they fit $B_R=1500$ tokens, and presented in temporal order. The
answerer returns an answer, the identifiers of the supporting records, and a
confidence; citations of records it was not shown are discarded and counted.
The time support of the cited records is the predicted evidence $I(C)$.
